% Source: https://artofproblemsolving.com/community/c3h1330670_arithmetic_sequence
% Stable identity: aops-topic-1330670-q01
Since $a+b+c+d=2(a+d)$, we want $d$ and $a$ as large as possible.

\textbf{A constant divisor-count sequence works.} The numbers $69=3\cdot23$, $77=7\cdot11$, $85=5\cdot17$, $93=3\cdot31$ form an arithmetic sequence with difference $8$, and each has $4$ divisors. This gives $a+b+c+d=324$.

\textbf{Nothing larger.} A sum above $324$ needs $a+d\ge163$; since $d\le99$ this forces $a\ge64$, so all four terms lie in $[64,99]$ and the common difference $e$ is at most $11$. There are only finitely many such sequences $a,a+e,a+2e,a+3e$, and checking their divisor counts (done exhaustively by computer) shows that none of them has $\tau$ values in arithmetic progression.

Hence the maximum is $\boxed{324}$. (If the divisor counts were required to form a \emph{non-constant} arithmetic sequence, the best would be $74,76,78,80$ with $4,6,8,10$ divisors, sum $308$.)
